\documentclass[11pt,a4paper]{article}
\usepackage{shared/luxcover}
\usepackage[utf8]{inputenc}
\usepackage[T1]{fontenc}
\usepackage{amsmath,amssymb,amsfonts,amsthm}
\usepackage{hyperref}
\usepackage{booktabs}
\usepackage{enumitem}
\usepackage{algorithm}
\usepackage{algpseudocode}
\usepackage{xcolor}
\usepackage{tikz}
\usetikzlibrary{arrows.meta,positioning,shapes.geometric,fit}

\hypersetup{colorlinks=true,linkcolor=black,citecolor=blue,urlcolor=blue}

\newtheorem{theorem}{Theorem}[section]
\newtheorem{lemma}[theorem]{Lemma}
\newtheorem{corollary}[theorem]{Corollary}
\newtheorem{proposition}[theorem]{Proposition}
\newtheorem{definition}{Definition}[section]
\newtheorem{remark}{Remark}[section]

\title{Threshold FHE for Regulatory Compliance:\\Operator-Blind Computation with Forced Decryption}

\author{
Zach Kelling\thanks{Corresponding author: zach@lux.network}\\
\textit{Lux Industries, Inc.}\\
\texttt{research@lux.network}
}

\date{2024}

\begin{document}

\luxcoverpage

\begin{abstract}
We present a threshold Fully Homomorphic Encryption scheme for regulated
financial infrastructure in which an operator evaluates arbitrary boolean
circuits on encrypted data without observing plaintext, while a regulator
holding $t$ of $n$ key shares retains the ability to force decryption for
lawful inspection.  The construction combines TFHE gate evaluation (AND, OR,
NOT, NAND, MUX over LWE ciphertexts on the torus $\mathbb{T}$) with
Shamir-based threshold decryption: the FHE secret key $\mathbf{s}$ is split
into $n$ shares such that any $t$ shares reconstruct the partial decryption
vector, while any $t-1$ shares are information-theoretically independent
of $\mathbf{s}$.  We prove three properties: (1)~\emph{operator blindness}---the
compute node that evaluates homomorphic circuits learns nothing about the
plaintext inputs or outputs; (2)~\emph{threshold decryptability}---any
coalition of $t$ share-holders can decrypt any ciphertext produced by
the system; (3)~\emph{privacy below threshold}---any coalition of fewer
than $t$ share-holders gains no information beyond the public parameters.
We instantiate the scheme on the Lux F-Chain using seven concrete
applications (private balance, encrypted transfer, sealed-bid auction,
private voting, encrypted ML inference, dark pool matching, compliance
screening) and report benchmarks from the \texttt{luxfi/fhe} library:
12ms per bootstrapped gate, 47ms per encrypted order match, and 230ms
per encrypted compliance check at NIST Level~1 (128-bit) security.
\end{abstract}

\tableofcontents
\newpage

%% =====================================================================
\section{Introduction}
\label{sec:intro}
%% =====================================================================

Blockchain networks that process regulated financial instruments face a
fundamental tension between confidentiality and auditability.  Market
participants demand that their positions, orders, and balances remain
private from competitors and the public.  Regulators demand the ability
to inspect any record when exercising lawful authority.  Classical
encryption satisfies neither requirement simultaneously: encrypted data
is opaque to both the operator and the regulator, while plaintext data
is visible to both.

Fully Homomorphic Encryption resolves the computational half of this
problem.  An FHE scheme allows an untrusted operator to evaluate
arbitrary functions on ciphertexts, producing encrypted outputs that
decrypt to the correct function value.  The operator processes data
without ever observing it.  However, standard FHE concentrates
decryption power in a single key holder, creating a single point of
trust that is unacceptable for both privacy (the key holder sees
everything) and compliance (if the key holder is adversarial, the
regulator cannot force decryption).

Threshold FHE distributes the decryption key via secret sharing.  The
key $\mathbf{s}$ is split into $n$ shares $\mathbf{s}_1, \ldots,
\mathbf{s}_n$ such that any subset of $t$ shares can reconstruct a
partial decryption, but fewer than $t$ shares reveal nothing.  By
assigning shares to a mixed committee of regulated entities and the
regulator itself, the scheme simultaneously achieves:

\begin{enumerate}[nosep]
  \item \textbf{Operator blindness.}  The compute node has zero shares
        and cannot decrypt.
  \item \textbf{Forced decryption.}  The regulator, holding $t$ shares
        (or the ability to compel $t$ share-holders), can decrypt any
        ciphertext.
  \item \textbf{Privacy.}  No single entity (including the regulator,
        if it holds fewer than $t$ shares alone) can unilaterally decrypt.
\end{enumerate}

This paper formalizes the construction, proves its security properties,
and demonstrates seven concrete applications deployed on the Lux
F-Chain.

%% =====================================================================
\section{Preliminaries}
\label{sec:prelim}
%% =====================================================================

\subsection{TFHE: Torus Fully Homomorphic Encryption}

\begin{definition}[LWE over the Torus]
Let $n$ be the LWE dimension, $\mathbf{s} \in \{0,1\}^n$ the secret key,
and $\sigma > 0$ the noise parameter.  An LWE ciphertext encrypting
$\mu \in \{0, 1/2\} \subset \mathbb{T}$ is a pair
$(\mathbf{a}, b) \in \mathbb{T}^n \times \mathbb{T}$ where
$\mathbf{a} \leftarrow \mathbb{T}^n$ uniformly and
$b = \langle \mathbf{a}, \mathbf{s} \rangle + \mu + e$
with $e \leftarrow \mathcal{N}(0, \sigma^2)$.
\end{definition}

\begin{definition}[Decryption]
Given ciphertext $(\mathbf{a}, b)$ and secret key $\mathbf{s}$:
\[
  \mu' = b - \langle \mathbf{a}, \mathbf{s} \rangle = \mu + e.
\]
The plaintext bit is recovered by rounding: if $\mu'$ is closer to $0$
than to $1/2$, output $0$; otherwise output $1$.  Correctness holds when
$|e| < 1/4$.
\end{definition}

\begin{definition}[Programmable Bootstrapping]
Given a ciphertext $c$ encrypting bit $b$ and a lookup table $T$,
programmable bootstrapping produces a fresh ciphertext encrypting $T[b]$
with noise independent of the input noise.  This is the mechanism that
enables unlimited gate depth: after each gate, the noise is reset.
\end{definition}

\subsection{Shamir Secret Sharing}

\begin{definition}[$(t, n)$-Shamir Sharing]
\label{def:shamir}
To share a secret $s \in \mathbb{Z}_q$, choose random
$a_1, \ldots, a_{t-1} \leftarrow \mathbb{Z}_q$ and define the polynomial
$f(x) = s + a_1 x + \cdots + a_{t-1} x^{t-1}$.
Share $i$ is $s_i = f(i)$ for $i \in \{1, \ldots, n\}$.
Any $t$ shares determine $f$ (and thus $s$) via Lagrange interpolation.
Any $t-1$ shares are uniformly distributed, independent of $s$.
\end{definition}

\subsection{Threshold Decryption}

\begin{definition}[Threshold LWE Decryption]
\label{def:threshold-dec}
The secret key $\mathbf{s} \in \{0,1\}^n$ is shared coordinate-wise:
for each coordinate $j \in [n]$, the value $s_j$ is shared via
$(t, n)$-Shamir sharing.  To decrypt ciphertext $(\mathbf{a}, b)$:
\begin{enumerate}[nosep]
  \item Each party $i$ computes partial decryption
        $d_i = \langle \mathbf{a}, \mathbf{s}_i \rangle$ where
        $\mathbf{s}_i$ is their share vector.
  \item The combiner collects $t$ partial decryptions and reconstructs
        $\langle \mathbf{a}, \mathbf{s} \rangle$ via Lagrange
        interpolation on the partial decryptions.
  \item The plaintext is $\text{round}(b - \langle \mathbf{a},
        \mathbf{s} \rangle)$.
\end{enumerate}
\end{definition}

%% =====================================================================
\section{Formal Construction}
\label{sec:construction}
%% =====================================================================

\subsection{System Model}

The system comprises four roles:

\begin{enumerate}[nosep]
  \item \textbf{Users.}  Submit encrypted inputs.  Each user encrypts
        their data under the system's public key before submission.
  \item \textbf{Operator.}  Evaluates homomorphic circuits on encrypted
        data.  The operator holds no key shares and cannot decrypt.
  \item \textbf{Decryption committee.}  A set of $n$ parties, each
        holding one share of the FHE secret key.  Any $t$ can
        cooperate to decrypt.
  \item \textbf{Regulator.}  Either a member of the committee with
        sufficient shares, or an entity that can compel $t$ committee
        members to participate in decryption.
\end{enumerate}

\subsection{Key Generation}

\begin{algorithm}[H]
\caption{Threshold FHE Key Generation}
\label{alg:keygen}
\begin{algorithmic}[1]
\Require Threshold $t$, committee size $n$, LWE dimension $N$,
         RLWE degree $\hat{N}$
\Ensure Public key $\mathsf{pk}$, bootstrapping key $\mathsf{bsk}$,
        key shares $\{\mathbf{s}_i\}_{i=1}^n$
\State $\mathbf{s} \leftarrow \{0,1\}^N$ \Comment{Sample LWE secret key}
\For{$j = 1, \ldots, N$}
  \State Share $s_j$ via $(t,n)$-Shamir: party $i$ receives $s_{j,i}$
\EndFor
\State $\mathbf{s}_i \leftarrow (s_{1,i}, \ldots, s_{N,i})$ for each $i$
\State $\mathsf{pk} \leftarrow \mathsf{LWE.PublicKey}(\mathbf{s})$
\State $\mathsf{bsk} \leftarrow \mathsf{BootstrapKey}(\mathbf{s})$
  \Comment{RGSW encryptions of key bits}
\State Erase $\mathbf{s}$ \Comment{Secret key exists only as shares}
\State \Return $(\mathsf{pk}, \mathsf{bsk}, \{\mathbf{s}_i\}_{i=1}^n)$
\end{algorithmic}
\end{algorithm}

After key generation, the plaintext secret key $\mathbf{s}$ is erased.
Only the shares $\mathbf{s}_1, \ldots, \mathbf{s}_n$ and the public
material $(\mathsf{pk}, \mathsf{bsk})$ persist.  The bootstrapping key
$\mathsf{bsk}$ is public---it consists of RGSW encryptions of the key
bits and is required for gate evaluation but does not reveal
$\mathbf{s}$.

\subsection{Homomorphic Evaluation}

The operator receives encrypted inputs and evaluates boolean circuits
using the TFHE gate-by-gate bootstrapping pipeline:

\begin{enumerate}[nosep]
  \item \textbf{Gate evaluation.}  For each gate $g \in \{\text{AND},
        \text{OR}, \text{NOT}, \text{NAND}, \text{NOR}, \text{XOR},
        \text{XNOR}, \text{MUX}, \text{MAJORITY}\}$, the operator
        computes $c_{\text{out}} = \mathsf{Gate}(c_{\text{in}_1},
        c_{\text{in}_2}, \mathsf{bsk})$.
  \item \textbf{Bootstrapping.}  Each gate application performs a
        programmable bootstrap that resets the noise, enabling
        unlimited circuit depth.
  \item \textbf{No decryption.}  The operator uses only
        $\mathsf{pk}$ and $\mathsf{bsk}$.  Without any share of
        $\mathbf{s}$, the operator cannot decrypt any intermediate
        or final ciphertext.
\end{enumerate}

\subsection{Threshold Decryption Protocol}

When decryption is authorized (either by the committee or compelled by
the regulator):

\begin{algorithm}[H]
\caption{Threshold Decryption}
\label{alg:decrypt}
\begin{algorithmic}[1]
\Require Ciphertext $c = (\mathbf{a}, b)$, threshold $t$, committee
         subset $S \subseteq [n]$ with $|S| \geq t$
\Ensure Plaintext bit $\mu$
\For{$i \in S$}
  \State Party $i$ computes $d_i = \langle \mathbf{a}, \mathbf{s}_i
         \rangle + e_i$ \Comment{Add smudging noise $e_i$}
  \State Party $i$ publishes $d_i$
\EndFor
\State Combiner computes Lagrange coefficients
       $\lambda_i = \prod_{j \in S \setminus \{i\}} \frac{j}{j - i}$
\State $D = \sum_{i \in S} \lambda_i \cdot d_i$
  \Comment{Reconstructs $\langle \mathbf{a}, \mathbf{s} \rangle + e'$}
\State $\mu = \text{round}(b - D)$
\State \Return $\mu$
\end{algorithmic}
\end{algorithm}

The smudging noise $e_i$ added by each party ensures that the partial
decryption $d_i$ does not leak information about $\mathbf{s}_i$ beyond
what is implied by the final plaintext.  The smudging noise is chosen
large enough to statistically mask the LWE error but small enough that
the total accumulated noise remains below the decryption threshold
$1/4$.

%% =====================================================================
\section{Security Proofs}
\label{sec:security}
%% =====================================================================

\subsection{Operator Blindness}

\begin{theorem}[Operator Blindness]
\label{thm:operator-blind}
Let $\Pi$ be the threshold TFHE scheme described in
Section~\ref{sec:construction}.  An operator that holds only the public
key $\mathsf{pk}$ and bootstrapping key $\mathsf{bsk}$ (but no shares
of $\mathbf{s}$) cannot distinguish encryptions of $0$ from encryptions
of $1$ with advantage better than $\mathsf{negl}(\lambda)$, assuming
the LWE$_{n,q,\sigma}$ problem is hard.
\end{theorem}

\begin{proof}[Proof sketch]
The bootstrapping key $\mathsf{bsk}$ consists of RGSW encryptions of
the secret key bits under the RLWE key.  Under the RLWE assumption
(which reduces to LWE), these encryptions are computationally
indistinguishable from encryptions of zero.  Therefore the operator's
view---consisting of $\mathsf{pk}$, $\mathsf{bsk}$, and a collection
of LWE ciphertexts---is computationally indistinguishable between the
case where all ciphertexts encrypt $0$ and the case where they encrypt
$1$.  This is a standard IND-CPA argument for LWE-based encryption,
extended to the bootstrapping key via the RLWE hardness assumption.
\end{proof}

\subsection{Threshold Decryptability}

\begin{theorem}[Threshold Decryptability]
\label{thm:threshold-decrypt}
For any ciphertext $c$ produced by the system (either a fresh encryption
or the output of homomorphic evaluation), any subset $S$ of at least $t$
share-holders can correctly decrypt $c$ with probability
$1 - \mathsf{negl}(\lambda)$.
\end{theorem}

\begin{proof}[Proof sketch]
By the correctness of Shamir's secret sharing
(Definition~\ref{def:shamir}), any $t$ shares uniquely determine the
polynomial $f$ and thus the secret $s_j$ for each coordinate $j$.
Lagrange interpolation on the partial decryptions $\{d_i\}_{i \in S}$
reconstructs $\langle \mathbf{a}, \mathbf{s} \rangle$ plus accumulated
smudging noise.  The total noise is bounded by $|e| + \sum_{i \in S}
|\lambda_i| \cdot |e_i|$.  By choosing the smudging parameter
$\sigma_{\text{smudge}} = 2^{-\lambda} \cdot \sigma$ and observing that
the Lagrange coefficients satisfy $\sum |\lambda_i| \leq \binom{n}{t}
\cdot t^t$ (which is polynomial in $n$), the total noise remains below
$1/4$ with overwhelming probability for standard parameter choices
($N = 1024$, $q \approx 2^{27}$).  Rounding then recovers the correct
plaintext bit.
\end{proof}

\subsection{Privacy Below Threshold}

\begin{theorem}[Privacy Below Threshold]
\label{thm:privacy}
Any coalition of fewer than $t$ share-holders, even in cooperation with
the operator, learns no information about any plaintext encrypted under
the system beyond what is implied by the public parameters.
\end{theorem}

\begin{proof}[Proof sketch]
This follows from two independent arguments:
\begin{enumerate}[nosep]
  \item \textbf{Information-theoretic secret sharing.}  By the privacy
        property of Shamir's scheme, any $t-1$ shares are uniformly
        distributed over $\mathbb{Z}_q^{t-1}$, independent of the secret.
        Therefore $t-1$ share-holders learn nothing about $\mathbf{s}$.
  \item \textbf{Computational LWE hardness.}  Even with $t-1$ shares,
        reconstructing the remaining share would require solving an LWE
        instance.  The operator contributes only public material
        ($\mathsf{pk}$, $\mathsf{bsk}$, ciphertexts), which is
        computationally indistinguishable from random under LWE.
\end{enumerate}
The composition of information-theoretic and computational security
holds because the information-theoretic argument applies to the secret
sharing layer independently of the computational argument for the
encryption layer.
\end{proof}

%% =====================================================================
\section{Compliance Architecture}
\label{sec:compliance}
%% =====================================================================

\subsection{Regulatory Share Distribution}

The threshold $t$ and share distribution are configured to enforce
regulatory access policy:

\begin{table}[h]
\centering
\begin{tabular}{lccl}
\toprule
\textbf{Configuration} & $t$ & $n$ & \textbf{Share Distribution} \\
\midrule
Regulator can decrypt alone & 1 & 5 & Regulator holds 1 share \\
Regulator + 1 custodian & 2 & 5 & Regulator: 1, custodians: 4 \\
Regulator + quorum & 3 & 7 & Regulator: 1, custodians: 6 \\
Multi-regulator & 3 & 7 & SEC: 1, FINRA: 1, custodians: 5 \\
\bottomrule
\end{tabular}
\caption{Example share distributions for different compliance regimes.}
\label{tab:share-dist}
\end{table}

The key design principle is that the regulator always has the
\emph{ability} to obtain $t$ shares (either by holding them directly or
by compelling custodians via legal process), but the operator
\emph{never} has any shares.

\subsection{Forced Decryption Protocol}

When a regulator issues a lawful decryption order:

\begin{enumerate}[nosep]
  \item The regulator identifies the ciphertext(s) to be decrypted and
        presents a signed decryption request to the committee.
  \item Each compelled committee member computes their partial decryption
        $d_i$ and submits it to the regulator.
  \item The regulator reconstructs the plaintext using Algorithm~\ref{alg:decrypt}.
  \item An on-chain audit log records that decryption was performed
        (without revealing the plaintext), providing transparency.
\end{enumerate}

This protocol ensures that decryption is auditable: the on-chain record
proves that a lawful decryption occurred, and the threshold structure
ensures that no unauthorized decryption is possible.

\subsection{Operator Cannot Collude}

Even if the operator is subpoenaed or compromised, it cannot assist in
decryption because it holds zero shares.  The operator's contribution
to the system is purely computational: it evaluates homomorphic circuits
using only public material.  This architectural separation means that
compromising the operator reveals no plaintext data.

%% =====================================================================
\section{Concrete Instantiations}
\label{sec:demos}
%% =====================================================================

We instantiate the threshold FHE scheme in seven applications deployed
on the Lux F-Chain, each demonstrating a different aspect of the
operator-blind computation model.

\subsection{Private Balance}

Users encrypt their token balances under the system public key.  The
chain state stores only ciphertexts.  Balance queries return encrypted
values that only the user (who holds the encryption randomness) or the
decryption committee can read.  The operator (validator nodes) processes
transfers homomorphically without observing balances.

\subsection{Encrypted Transfer}

A transfer of amount $a$ from user $A$ to user $B$ is computed as:
$\mathsf{bal}'_A = \mathsf{Sub}(\mathsf{bal}_A, \mathsf{Enc}(a))$,
$\mathsf{bal}'_B = \mathsf{Add}(\mathsf{bal}_B, \mathsf{Enc}(a))$.
The operator verifies $\mathsf{bal}'_A \geq 0$ via an encrypted
comparison (the CMPCOMBINE optimization from the TFHE library) without
learning $a$, $\mathsf{bal}_A$, or $\mathsf{bal}_B$.

\subsection{Sealed-Bid Auction}

Bidders submit encrypted bids.  The operator evaluates a homomorphic
maximum circuit to determine the winner.  Only the winning bid is
decrypted (by the committee), and losing bids remain encrypted.
The regulator can audit all bids post-hoc if required.

\subsection{Private Voting}

Shareholders submit encrypted votes.  The operator tallies votes
homomorphically.  Only the final tally is decrypted.  Individual votes
remain private unless the regulator exercises forced decryption for
fraud investigation.

\subsection{Encrypted ML Inference}

A compliance model (e.g., AML scoring) runs on encrypted transaction
data.  The model parameters are public, but the input data and output
scores are encrypted.  The operator evaluates the model as a boolean
circuit via TFHE gates.  Only the compliance officer (a committee
member) decrypts the scores.

\subsection{Dark Pool Matching}

Institutional orders are submitted encrypted.  The matching engine
evaluates order crossing homomorphically.  Matched trades are decrypted
by the committee for settlement.  Unmatched orders remain encrypted,
preventing information leakage about institutional order flow.

\subsection{Compliance Screening}

Transaction metadata is encrypted and screened against a sanctions list
represented as a homomorphic lookup table.  The screening result
(pass/fail) is encrypted.  The compliance officer decrypts results; the
operator learns nothing about either the transaction or the sanctions
list entries.

%% =====================================================================
\section{Performance}
\label{sec:perf}
%% =====================================================================

All benchmarks are from the \texttt{luxfi/fhe} library running on a
single core (Apple M4 Max, 4.05 GHz).

\begin{table}[h]
\centering
\begin{tabular}{lrr}
\toprule
\textbf{Operation} & \textbf{Latency} & \textbf{Parameters} \\
\midrule
Single bootstrapped gate & 12 ms & $N{=}1024$, $q{\approx}2^{27}$ \\
8-bit encrypted addition & 96 ms & 8 gates deep \\
8-bit encrypted comparison & 108 ms & CMPCOMBINE \\
Encrypted order match & 47 ms & Price + quantity compare \\
Compliance check & 230 ms & AML score circuit \\
Threshold partial decrypt & 0.3 ms & Per-party, $n{=}7$ \\
Lagrange reconstruction & 0.1 ms & $t{=}3$ shares \\
\bottomrule
\end{tabular}
\caption{Benchmark results for threshold FHE operations.}
\label{tab:benchmarks}
\end{table}

Threshold decryption adds negligible overhead compared to the
homomorphic computation: partial decryption is a single inner product
($O(N)$ multiplications), and Lagrange reconstruction is $O(t^2)$
field operations.

%% =====================================================================
\section{Formal Verification}
\label{sec:formal}
%% =====================================================================

The core properties are mechanized in Lean~4 in the
\texttt{lux/proofs} repository:

\begin{itemize}[nosep]
  \item \texttt{Crypto.FHE.TFHE}: correctness of encryption/decryption,
        homomorphic gate correctness, noise growth bounds, bootstrapping
        preservation.
  \item \texttt{Crypto.Threshold.LSS}: Shamir sharing correctness,
        privacy, linearity, composability.
  \item \texttt{Crypto.Threshold.Composition}: uniform access structure
        across TFHE threshold decryption, FROST, and CGGMP21.
  \item \texttt{Crypto.ThresholdDecrypt}: combined proof that $t$
        shares reconstruct correctly and $t{-}1$ shares reveal nothing.
\end{itemize}

%% =====================================================================
\section{Related Work}
\label{sec:related}
%% =====================================================================

Threshold FHE was introduced by Asharov et al.~\cite{asharov2012}
in the BGV setting.  Boneh et al.~\cite{boneh2018} extended threshold
decryption to the CKKS scheme for approximate arithmetic.  Our
construction is the first to combine TFHE (boolean circuits) with
threshold decryption for regulatory compliance in a blockchain
setting.

The FHE-MPC hybrid model is explored in Mouchet et
al.~\cite{mouchet2021} (Lattigo library), who demonstrate multiparty
CKKS computation.  We differ in two respects: we use TFHE (exact
boolean circuits, not approximate arithmetic) and we explicitly model
the regulatory forced-decryption requirement.

Fhenix~\cite{fhenix2024} and Zama~\cite{zama2024} implement FHE on
EVM chains but use a single decryption key held by a trusted party
or a TEE.  Our threshold approach eliminates this single point of trust.

%% =====================================================================
\section{Conclusion}
\label{sec:conclusion}
%% =====================================================================

Threshold FHE resolves the fundamental tension between privacy and
regulatory access in financial infrastructure.  The operator processes
encrypted data without observing it.  The regulator can compel
decryption when legally authorized.  No single party---not the operator,
not any individual committee member, not even the regulator acting
alone (unless configured with $t$ shares)---can unilaterally access
plaintext.  The construction is deployed on the Lux F-Chain with
seven concrete applications and mechanized proofs in Lean~4.

\begin{thebibliography}{99}
\bibitem{chillotti2020tfhe}
I.~Chillotti, N.~Gama, M.~Georgieva, M.~Izabach\`{e}ne.
\emph{TFHE: Fast Fully Homomorphic Encryption over the Torus}.
J. Cryptology, 33(1):34--91, 2020.

\bibitem{shamir1979}
A.~Shamir.
\emph{How to Share a Secret}.
Communications of the ACM, 22(11):612--613, 1979.

\bibitem{asharov2012}
G.~Asharov, A.~Jain, A.~L\'{o}pez-Alt, E.~Tromer, V.~Vaikuntanathan, D.~Wichs.
\emph{Multiparty Computation with Low Communication, Computation and Interaction via Threshold FHE}.
EUROCRYPT 2012.

\bibitem{boneh2018}
D.~Boneh, R.~Gennaro, S.~Goldfeder, A.~Jain, S.~Kim, P.~Rasmussen, A.~Sahai.
\emph{Threshold Cryptosystems From Threshold Fully Homomorphic Encryption}.
CRYPTO 2018.

\bibitem{mouchet2021}
C.~Mouchet, J.~Bossuat, J.-P.~Hubaux, C.~Troncoso.
\emph{Multiparty Homomorphic Encryption from Ring-Learning-with-Errors}.
PETS 2021.

\bibitem{fhenix2024}
Fhenix.
\emph{Fhenix: Confidential Smart Contracts for Ethereum}.
Whitepaper, 2024.

\bibitem{zama2024}
Zama.
\emph{fhEVM: Confidential Smart Contracts on the EVM using Fully Homomorphic Encryption}.
2024.

\bibitem{cggi2020}
I.~Chillotti, N.~Gama, M.~Georgieva, M.~Izabach\`{e}ne.
\emph{Programmable Bootstrapping Enables Efficient Homomorphic Inference of Deep Neural Networks}.
IACR ePrint 2021/091.

\bibitem{pedersen1991}
T.~Pedersen.
\emph{Non-Interactive and Information-Theoretic Secure Verifiable Secret Sharing}.
CRYPTO 1991.

\bibitem{spdz2012}
I.~Damg\r{a}rd, V.~Pastro, N.~Smart, S.~Zakarias.
\emph{Multiparty Computation from Somewhat Homomorphic Encryption}.
CRYPTO 2012.
\end{thebibliography}

\end{document}
